% The clearance horizon H: a revised derivation (T024).
% Build: tectonic horizon_derivation.tex   (or: latexmk -pdf horizon_derivation.tex)
% Every number in this document is printed by one of two scripts:
%   uv run python scripts/spike_t024_object_motion.py   (data evidence, Sections 2 and 5)
%   uv run python scripts/t024_worked_examples.py       (worked examples, Section 8)
\documentclass[11pt]{article}
\usepackage[a4paper,margin=2.4cm]{geometry}
\usepackage{amsmath,amssymb,amsthm,mathtools}
\usepackage{booktabs}
\usepackage{enumitem}
\usepackage{xcolor}
\usepackage[colorlinks=true,linkcolor=blue!50!black,citecolor=blue!50!black,urlcolor=blue!50!black]{hyperref}

\newtheorem{lemma}{Lemma}
\newtheorem{proposition}{Proposition}
\newtheorem{theorem}{Theorem}
\theoremstyle{definition}
\newtheorem{definition}{Definition}
\newtheorem{assumption}{Assumption}
\theoremstyle{remark}
\newtheorem{remark}{Remark}

\newcommand{\R}{\mathbb{R}}
\newcommand{\Disk}{\mathcal{D}}
\newcommand{\Occ}{\mathcal{O}}
\newcommand{\Ohat}{\widehat{\mathcal{O}}}
\newcommand{\Ego}{\mathcal{E}}
\newcommand{\Hc}{H_{\mathrm{c}}}
\newcommand{\Hr}{H_{\mathrm{r}}}
\newcommand{\Tcap}{T_{\mathrm{cap}}}
\newcommand{\pref}{p_{\mathrm{ref}}}
\newcommand{\psiref}{\psi_{\mathrm{ref}}}
\newcommand{\vref}{v_{\mathrm{ref}}}
\DeclareMathOperator{\dist}{dist}
\DeclareMathOperator{\Rect}{Rect}

\title{The clearance horizon $H$:\\[2pt] a revised derivation}
\author{T024 --- draft for review}
\date{12 November 2026}

\begin{document}
\maketitle

\begin{abstract}
When the car receives a plan from the Alpamayo server, it computes once a horizon $H$: how long the
plan may be followed before the next plan must have arrived. The baseline equations
(\texttt{Deadline\_Aware\_Multi\_Vehicle\_Serving.pdf}) give $H=0$ for 52\% of the 1\,005 plans of our
five clips. This document explains why, using our data, and derives a replacement. The ego and the
objects become oriented boxes instead of disks; the object bounds become per-axis, conditioned on the
object's class and motion, and calibrated on data with a stated probability; and the horizon is
split into a \emph{contact horizon} $\Hc$ (when could something first touch the plan?) and a
\emph{reaction horizon} $\Hr$ (until when can the car still avoid it?), of which only $\Hr$ is a safe
deadline. Every step is proved or measured; every number is reproduced by a script in the repository.
\end{abstract}

\tableofcontents

%======================================================================
\section*{At a glance}
\addcontentsline{toc}{section}{At a glance}

\paragraph{What goes wrong in the baseline.} Measured on the recorded stand-in (5 clips, offline ego
poses, 1\,005 plans), using the T014 runner:
\begin{enumerate}[leftmargin=*]
  \item \textbf{Geometry, not motion, makes $H=0$.} For 516 plans $H=0$. In 98.4\% of them the real boxes
    of the ego and of the limiting object are more than 0.5\,m apart at $\tau=0$ (median 1.48\,m). The
    disks need $6.4$\,m between centres; a car in the next lane is about $3.5$\,m away.
  \item \textbf{The object bound has the wrong shape.} Prediction errors grow about linearly in $\tau$
    (velocity error dominates), not quadratically. The baseline $\beta$ is too small up to 1\,s (1.7--1.8\%
    of vehicle cases fall outside it) and far too large later (44\,m at 6.4\,s). It is the same in every
    direction, although a stopped car's box drifts about 5 times less than a moving car's.
  \item \textbf{The ego allowance grows linearly forever}, to 5.3\,m at 6.4\,s at 10\,m/s, although a
    tracking controller keeps the error bounded.
  \item \textbf{``First possible contact'' leaves no time to act.} If the next plan arrives exactly at
    $t+H$, a contact may already be unavoidable.
  \item \textbf{No finite worst-case bound holds} on heavy-tailed data; the baseline's guarantee is a
    probability statement in disguise, with an unknown probability.
\end{enumerate}

\paragraph{What the revision does.}
\begin{enumerate}[leftmargin=*,label=(\alph*)]
  \item Oriented boxes for ego and objects (the geometry the safety check T015 already uses). The ego's
    tracking error enters as $\rho_p + 2R_F\sin(\bar e_\theta/2) = 0.39$\,m: the 3.79\,m radius moves from
    the body size into the heading-error term (Lemma~\ref{lem:ego}).
  \item Each object's predicted occupancy is a box in its own frame with separate front, back and side
    allowances. It is the intersection of a \emph{physical} layer (what the object can do) and a
    \emph{calibrated} layer (what objects of its kind did in our data), the latter with a
    distribution-free coverage guarantee (conformal prediction).
  \item Two horizons: $\Hc$, and $\Hr\le\Hc$, the last moment from which braking along the plan still
    avoids every predicted occupancy. The scheduler's deadline should use $\Hr$.
  \item Explicit assumptions (who is responsible for which contact) and an explicit guarantee:
    $\Pr(\text{contact with a known object}) \le \sum_k \delta_k$.
\end{enumerate}

\paragraph{Effect on four example scenes} (Section~\ref{sec:examples}; $H$ in seconds):

\begin{center}
\begin{tabular}{p{6.4cm}cccc}
\toprule
Scene & baseline & physical & calibrated p99 & both \\
\midrule
E1 car in the next lane, abreast, 10\,m/s & 0.00 & 0.09 & 0.60 & 0.60 \\
E2 following at 4\,m/s, 4\,m gap & 1.36 & 0.50 & 1.85 & 1.85 \\
\quad reaction horizon $\Hr$ & --- & 0.03 & 1.66 & 1.66 \\
E3 pedestrian crossing (hand example) & 2.00 & 2.23 & 2.57 & 2.57 \\
E4 passing a parked car at 10\,m/s & 0.76 & 0.75 & 6.40 (capped) & 6.40 (capped) \\
\bottomrule
\end{tabular}
\end{center}
The calibrated numbers are provisional: fitted on the same five clips they describe, at a per-$\tau$
level rather than the joint level of Section~\ref{sec:conformal}.

%======================================================================
\section{Setting and notation}\label{sec:setting}

\subsection{The system}
At time $t$ the car receives a plan from the server. It computes $H$ once, from what it knows at $t$,
and attaches the deadline $t+H$ to its next inference request; the server's scheduler serves the
earliest deadline first (T023). Two properties are wanted:
\begin{description}[leftmargin=*]
  \item[(R1) Safe:] if the next plan arrives by $t+H$, the car can still avoid contact with every object
    it knew about at $t$ --- except with a small, stated probability.
  \item[(R2) Useful:] subject to (R1), $H$ is as large as possible.
\end{description}
$H$ may only use data known at $t$ (the causality rule). The \emph{parameters} of the equations may be
calibrated offline on other data, like any trained model; to keep the evaluation honest they must be
calibrated on clips that are not used to evaluate $H$.

\subsection{What we know at $t$}
\begin{center}
\begin{tabular}{p{4.2cm}p{6.0cm}p{4.6cm}}
\toprule
Symbol & Meaning & Source \\
\midrule
$\pref(\tau),\ \psiref(\tau),\ \vref(\tau)$ & plan: rear-axle position, heading, speed, $\tau\in[0,T_p]$ & reference (T011/T012) \\
$F_e=[x_-,x_+]\times[y_-,y_+]$ & ego footprint about the rear axle & scene \texttt{meta.json} \\
$R_F=\max_{r\in F_e}\|r\|$ & farthest footprint point from the rear axle & $=3.790$\,m \\
$\hat c_k,\ \hat\psi_k$ & object box centre and heading at $t$ & \texttt{objects.parquet} $x,y,\mathit{yaw}$ \\
$\ell_k,\ w_k$ & half-length and half-width of the box & \texttt{length}/2, \texttt{width}/2 \\
$\hat w_k,\ \hat v_k=\hat w_k\cdot u_k$ & velocity; speed along the heading & $v_x,v_y$ (backward difference, T009) \\
$u_k,\ n_k$ & unit vectors along and across $\hat\psi_k$ & $u_k=(\cos\hat\psi_k,\sin\hat\psi_k)$ \\
class, motion state & vehicle / person / \dots; stopped if $\|\hat w_k\|<0.5$\,m/s & \texttt{class\_group}, \texttt{label\_class} \\
$\mathcal K(t)$ & known objects: \texttt{in\_scene} at $t$, velocity known & T008, T014 \\
\bottomrule
\end{tabular}
\end{center}
Our ego footprint is $x\in[-1.032,\,3.656]$, $y\in[-0.9995,\,0.9995]$\,m (4.69\,m $\times$ 2.00\,m).
Alpamayo gives a rotation per step, so $\psiref$ is known; for the recorded stand-in it is the scene's
yaw.

\paragraph{Sets.} $\Disk(c,r)$ is the closed disk of centre $c$ and radius $r$; $A\oplus B=\{a+b\}$ is
the Minkowski sum; $\Disk_r=\Disk(0,r)$; $\dist(A,B)=\min_{a\in A,b\in B}\|a-b\|$.
$\Rect(c,\psi,a,b)$ is the rectangle of centre $c$, heading $\psi$, half-length $a$ and half-width $b$.
$\Occ_k(\tau)$ is the region object $k$ really occupies at $t+\tau$, and $\Ohat_k(\tau)$ the region we
predict for it. $\Ego(\tau)$ is the region the ego is guaranteed to stay in.

%======================================================================
\section{The baseline, and what goes wrong}\label{sec:baseline}

\subsection{The baseline in one line}
In the notation of the PDF, for every known object,
\begin{equation}
m_k(\tau)=\big\|\pref(\tau)-(\hat c_k+\hat w_k\tau)\big\|-\Big[r_{\mathrm{veh}}+r_k+\rho_0+\varepsilon_p+(\bar q+\varepsilon_v)\tau+\tfrac12\bar a\tau^2\Big],
\qquad H=\inf\{\tau\le\Tcap:\min_k m_k(\tau)\le0\},
\label{eq:baseline}
\end{equation}
with $\bar q=\bar e_v+2\bar v_{\mathrm{ref}}\sin(\bar e_\theta/2)$, $r_{\mathrm{veh}}$ the disk about the
rear axle that covers the footprint ($3.790$\,m) and $r_k$ half the object's diagonal.

\subsection{What is right, and kept}
\begin{itemize}[leftmargin=*]
  \item The kinematic (no-slip) model of the ego and the error variables $e_p,e_\theta,e_v$.
  \item The velocity-mismatch bound $\|\dot p-\dot p_{\mathrm{ref}}\|\le|e_v|+2\vref|\sin(e_\theta/2)|$:
    correct, and the chord identity $\|h(\theta)-h(\theta_{\mathrm{ref}})\|=2|\sin(e_\theta/2)|$ behind
    it is reused in Lemma~\ref{lem:ego}.
  \item The Taylor-with-remainder bound
    $\|o(t+\tau)-\hat o-\hat w\tau\|\le\varepsilon_p+\varepsilon_v\tau+\tfrac12\bar a\tau^2$: correct for an
    object that can accelerate in any direction. It is kept for pedestrians
    (Proposition~\ref{prop:ped}).
  \item $H$ as an infimum over $[0,\Tcap]$, $\Tcap=\min(T_{\mathrm{rem}},T_{\mathrm{valid}})$.
\end{itemize}

\subsection{Problem 1: disks make geometry, not motion, the limit}
The ego disk is centred on the rear axle, so it must reach the front corners: $r_{\mathrm{veh}}=3.79$\,m.
Its area is $45.1\,\text{m}^2$, 4.8 times the footprint ($9.4\,\text{m}^2$); sideways it extends 3.79\,m
where the car extends 1.00\,m. A median car in our data is $3.89\times1.77$\,m, so $r_k=2.14$\,m (2.1
times the box area). At $\tau=0$ the baseline requires
$3.79+2.14+0.2+0.3=6.43$\,m between the ego's rear axle and the car's centre; the centre of a car in
the next lane is about 3.5\,m to the side.

\textbf{Data.} Of the 516 plans with $H=0$, 461 are limited by an object beside the ego
($|y_{\mathrm{rel}}|\ge1.5$\,m), 431 by a vehicle and 85 by a pedestrian or rider. The real box gap at $\tau=0$
between the ego and the limiting object is: p10 0.95\,m, median 1.48\,m, p90 3.67\,m; only 1.6\% of
them actually touch (\texttt{spike\_t024\_object\_motion.py}, part 5).

\subsection{Problem 2: the object bound has the wrong shape}
Write the constant-velocity residual $r_k(\tau)=c_k(t+\tau)-\hat c_k-\hat w_k\tau$. Table~\ref{tab:beta}
compares $\|r_k\|$ for every known vehicle at every tick of the five clips with the baseline
$\beta(\tau)=0.3+0.5\tau+\tau^2$.

\begin{table}[h]
\centering
\begin{tabular}{rrrrrr}
\toprule
$\tau$ [s] & $\beta(\tau)$ [m] & p99 $\|r\|$ & p99.9 $\|r\|$ & share above $\beta$ & samples (tracks) \\
\midrule
0.5 & 0.80 & 1.05 & 2.45 & 1.80\% & 55\,515 (692) \\
1.0 & 1.80 & 2.22 & 5.06 & 1.66\% & 52\,616 (628) \\
2.0 & 5.30 & 4.24 & 8.86 & 0.48\% & 47\,317 (568) \\
4.0 & 18.30 & 9.59 & 18.37 & 0.10\% & 37\,909 (468) \\
6.4 & 44.46 & 19.17 & 34.38 & 0.02\% & 28\,447 (381) \\
\bottomrule
\end{tabular}
\caption{Vehicles: constant-velocity residual against the baseline $\beta$. Samples from one track are
strongly correlated; the track count is the honest sample size.}
\label{tab:beta}
\end{table}

Fitting $c_0+c_1\tau+c_2\tau^2$ ($c_i\ge0$) to the quantiles gives
p99 $\approx0.60+1.21\tau+0.26\tau^2$ and p99.9 $\approx1.59+2.70\tau+0.37\tau^2$, against the
baseline's $0.3+0.5\tau+1.0\tau^2$. In the baseline's terms the data says $\varepsilon_v\approx1.2$--$2.7$
m/s (not 0.5) and $\bar a\approx0.5$--$0.7$\,m/s$^2$ (not 2): \emph{velocity error dominates}. T009 already
measured object velocity noise at 0.47--0.72\,m/s p95 from parked tracks, and showed that it is slow
drift that averaging does not remove.

Three more mismatches:
\begin{itemize}[leftmargin=*]
  \item \textbf{$\bar a$ is not a bound on braking.} Over 1\,s windows, vehicle speed changes reach
    $-4.67$\,m/s$^2$ at p0.1 and $-7.20$\,m/s$^2$ at the minimum, and cars with ABS brake at
    8.5--10.6\,m/s$^2$ on dry roads \cite{miller2023}. A worst-case bound with $\bar a=10$ would give
    $\beta(2\,\text{s})=21.3$\,m.
  \item \textbf{Isotropy.} At p90 the lateral residual (across the heading) is
    $0.09+0.13\tau+0.04\tau^2$ and the longitudinal one $0.41\tau+0.04\tau^2$: three times larger along
    the heading. A parked car cannot move sideways or backwards at all.
  \item \textbf{One size for all objects.} At $\tau=1$\,s, p99, the true box of a \emph{stopped}
    vehicle sticks out of its predicted box by 0.43\,m sideways; a \emph{moving} vehicle's by 2.02\,m
    (Table~\ref{tab:overhang}).
\end{itemize}

\subsection{Problem 3: the ego allowance grows linearly forever}
The PDF integrates the velocity mismatch open-loop: $\|e_p(\tau)\|\le\rho_0+\bar q\tau$. With
$\bar e_\theta=0.05$\,rad at 10\,m/s, $\bar q=0.8$\,m/s and $\rho(6.4)=5.3$\,m. The bound is correct
\emph{if} the heading and speed errors stay within $\bar e_\theta,\bar e_v$ and nothing corrects them;
but a tracking controller is precisely what keeps $e_p$ small. Closed-loop, the standard object is a
bounded \emph{tube}: $\|e_p\|\le\bar\rho_p$ for all $\tau$ \cite{mayne2005}. Note also that on our
dataset the ego follows the plan exactly, so this term is never tested (T018).

\subsection{Problem 4: ``first possible contact'' leaves no time to act}
The baseline's $H$ is the first time at which a contact \emph{becomes possible}. Suppose the next plan
arrives exactly at $t+H$. At that instant the ego's region already touches an object's region, whatever
the new plan says: a car cannot leave its current position instantly, and at 10\,m/s it needs
$v^2/(2b)=8.3$\,m and 1.7\,s to stop at $b=6$\,m/s$^2$. So a deadline equal to the first possible contact
violates (R1). What is needed is the \emph{time-to-react}: the latest moment from which an evasive
manoeuvre still succeeds \cite{hillenbrand2006,sontges2018}. Section~\ref{sec:horizons} defines it
for our setting.

\subsection{Problem 5: worst-case bounds and heavy tails}
The largest vehicle residual at 6.4\,s is 95\,m (a labelling artefact), and at 1\,s it is 14\,m. No
finite $\varepsilon_v,\bar a$ covers every case, so the baseline's ``guarantee'' really holds with some
probability, but that probability is neither chosen nor known. The revision states it: the bounds are
calibrated to a chosen miss probability $\delta$, with a guarantee that does not depend on the shape of
the error distribution (Section~\ref{sec:conformal}). Physics still provides hard limits where they
exist (a parked car does not drive backwards at 10\,m/s).

\subsection{Problem 6: two geometries for one question}
$H$ uses disks while the safety check (T015) uses boxes. The same geometry should be used for both, so
that a failure means ``a prediction was wrong'' and never ``the shapes differ''.

%======================================================================
\section{Geometry: boxes instead of disks}\label{sec:geometry}

\subsection{The ego's region}
\begin{lemma}[Ego body under tracking error]\label{lem:ego}
Let the true rear-axle pose be $(p,\theta)$ and the plan's $(\pref,\psiref)$, with
$\|p-\pref\|\le\rho_p$ and $|\theta-\psiref|\le\bar e_\theta\le\pi$. Then every point of the true body
lies in
\[
\Rect_{\mathrm{ref}}\oplus\Disk_{\rho},\qquad \rho=\rho_p+2R_F\sin(\bar e_\theta/2),
\]
where $\Rect_{\mathrm{ref}}$ is the footprint placed at the plan's pose.
\end{lemma}
\begin{proof}
A body point with coordinates $r\in F_e$ is at $x=p+R(\theta)r$; its planned position is
$x_{\mathrm{ref}}=\pref+R(\psiref)r\in\Rect_{\mathrm{ref}}$. Then
$x-x_{\mathrm{ref}}=(p-\pref)+R(\psiref)\big(R(e_\theta)-I\big)r$ with $e_\theta=\theta-\psiref$. A rotation
by $e_\theta$ moves a vector of length $\|r\|$ by the chord $2|\sin(e_\theta/2)|\,\|r\|$ (the PDF's
identity), and $\|r\|\le R_F$. Hence $\|x-x_{\mathrm{ref}}\|\le\rho_p+2R_F\sin(\bar e_\theta/2)$.
\end{proof}
With $\rho_p=0.2$\,m and $\bar e_\theta=0.05$\,rad: $\rho=0.2+2\cdot3.790\cdot\sin(0.025)=0.2+0.189=0.389$\,m.
The 3.79\,m appears only multiplied by $2\sin(\bar e_\theta/2)\approx0.05$.

\paragraph{Which $\rho_p$?} Closed-loop, $\rho_p$ is the controller's tube radius, a constant; it should
come from the controller's specification and be validated when a controller exists (T018). The PDF's
open-loop bound $\rho_p(\tau)=\rho_0+\bar q\tau$ remains available as a conservative option when
nothing is known about the controller. In the dataset replay the ego follows the plan exactly, so
$\rho$ only makes $H$ more conservative; the safety check uses $\rho=0$ (T023).

\subsection{Separation of inflated sets}
\begin{lemma}\label{lem:inflate}
For compact $A,B\subset\R^2$ and $a,b\ge0$:
$(A\oplus\Disk_a)\cap(B\oplus\Disk_b)=\emptyset\iff\dist(A,B)>a+b$.
\end{lemma}
\begin{proof}
The sets intersect iff there are $x\in A$, $y\in B$, $\|u\|\le a$, $\|v\|\le b$ with $x+u=y+v$, i.e.\ iff
some $\|x-y\|\le a+b$; by compactness the minimum $\dist(A,B)$ is attained.
\end{proof}
Moreover $\Rect(c,\psi,a,b)\oplus\Disk_\rho\subseteq\Rect(c,\psi,a+\rho,b+\rho)$: the enlarged rectangle
only adds the four corner squares, $(4-\pi)\rho^2=0.13\,\text{m}^2$ for $\rho=0.389$\,m. We use the
enlarged rectangle, so that every test below is between two rectangles, or a rectangle and a disk.

\subsection{Testing two rectangles}
Two convex polygons are disjoint iff some edge normal of one of them separates their projections
(the separating-axis theorem, used for 3-D boxes with 15 candidate axes in \cite{gottschalk1996}; for
two rectangles in the plane, their 4 edge normals suffice). When they are disjoint,
their distance is the smallest distance from a vertex of one to an edge of the other (16 + 16 pairs).
General convex sets: GJK \cite{gilbert1988}. A rectangle against a disk $\Disk(c,r)$: they are disjoint
iff the distance from $c$ to the rectangle (clamp $c$ in the rectangle's frame) exceeds $r$. All of this
is constant work per pair and per time step.

An alternative that keeps disks: following \cite{ziegler2010}, cover the footprint with $n$ equal disks
along its length; elementary geometry gives the radius $\sqrt{(L/2n)^2+(W/2)^2}$ (also Eq.~4 of
\cite{tolksdorf2024}). For our ego ($L=4.69$, $W=2.00$): 2.55\,m ($n=1$, centred on
the box, not the rear axle), 1.54\,m ($n=2$), 1.27\,m ($n=3$). Even $n=3$ loses about 0.5\,m against the
exact boxes beside a car in the next lane, and the safety check needs exact boxes anyway; we use boxes.

\subsection{A rotating box}
\begin{lemma}[Rotation envelope]\label{lem:rot}
Let a rectangle of half-dimensions $\ell\ge w$ be rotated about its centre by any angle
$|\phi|\le\bar\phi\le\pi/2$. All such rectangles lie in the rectangle with the original axes and
half-dimensions
\[
\ell'(\bar\phi)=\begin{cases}\ell\cos\bar\phi+w\sin\bar\phi & \bar\phi<\arctan(w/\ell)\\ \sqrt{\ell^2+w^2} & \text{otherwise}\end{cases},
\qquad
w'(\bar\phi)=\begin{cases}w\cos\bar\phi+\ell\sin\bar\phi & \bar\phi<\arctan(\ell/w)\\ \sqrt{\ell^2+w^2} & \text{otherwise.}\end{cases}
\]
\end{lemma}
\begin{proof}
The rotated rectangle's extent along the original long axis is
$h_u(\phi)=\ell|\cos\phi|+w|\sin\phi|=\sqrt{\ell^2+w^2}\cos(|\phi|-\arctan(w/\ell))$, which increases up to
$|\phi|=\arctan(w/\ell)$ and decreases after; across it,
$h_n(\phi)=w|\cos\phi|+\ell|\sin\phi|=\sqrt{\ell^2+w^2}\sin(|\phi|+\arctan(w/\ell))$, which increases up to
$|\phi|=\arctan(\ell/w)$. A set whose projections on both axes are bounded lies in the corresponding
rectangle.
\end{proof}

%======================================================================
\section{The objects: predicted occupancy}\label{sec:objects}

Every predicted region below is a rectangle with the object's axes $(u_k,n_k)$ at $t$, described by a
longitudinal interval $[L_k^-(\tau),L_k^+(\tau)]$ (along $u_k$, measured from $\hat c_k$) and a lateral
half-extent $\Lambda_k(\tau)$:
\[
\Ohat_k(\tau)=\Big\{\hat c_k+\alpha u_k+\beta n_k:\ \alpha\in[L_k^-,L_k^+],\ |\beta|\le\Lambda_k\Big\}.
\]
Two such rectangles with the same axes intersect in a rectangle of the same form, with the intervals
intersected. This is how the two layers are combined (Section~\ref{sec:combine}).

\subsection{Physical layer: vehicles}
\begin{assumption}[Vehicles]\label{ass:veh}
For a known vehicle $k$, over $[0,T]$:
\begin{enumerate}[label=(V\arabic*),leftmargin=*]
  \item \emph{Estimates at $t$:} $\|c_k(t)-\hat c_k\|\le\varepsilon_p$, $|\psi_k(t)-\hat\psi_k|\le\varepsilon_\psi$,
    $|v_k(t)-\hat v_k|\le\varepsilon_v$.
  \item \emph{No slip, forward:} the rear-axle point $P_k$ moves with velocity $v_k\,(\cos\psi_k,\sin\psi_k)$,
    $v_k\ge0$ (no reversing within the horizon; constraint C3 of \cite{althoff2016}).
  \item \emph{Longitudinal limits:} $-b_{\max}\le\dot v_k\le a_{\max}$ and $v_k\le v_{\max}$ (C4 and C1 of
    \cite{althoff2016}, who also limit acceleration by engine power and keep vehicles on the road).
  \item \emph{Turning limits:} curvature $|\kappa_k|\le\kappa_{\max}$ and lateral acceleration
    $v_k^2|\kappa_k|\le a_\perp$, i.e.\ $|\dot\psi_k|\le\omega(v_k)=\min(\kappa_{\max}v_k,\ a_\perp/v_k)$.
  \item \emph{Box:} the box is rigid, with its centre $c_k=P_k+l_k(\cos\psi_k,\sin\psi_k)$, $0\le l_k\le\ell_k$.
\end{enumerate}
\end{assumption}

Define, for $s\in[0,T]$,
\begin{align*}
v^+(s)&=\min\big(v_{\max},\ \hat v_k+\varepsilon_v+a_{\max}s\big), & v^-(s)&=\max\big(0,\ \hat v_k-\varepsilon_v-b_{\max}s\big),\\
\bar\omega(s)&=\max_{v\in[v^-(s),\,v^+(s)]}\omega(v), & \Omega(\tau)&=\int_0^\tau\bar\omega(s)\,ds,\qquad \bar\Delta(\tau)=\varepsilon_\psi+\Omega(\tau),\\
S^+(\tau)&=\int_0^\tau v^+(s)\,ds, & S^-(\tau)&=\int_0^\tau \underline{g}(s)\,ds,\quad
\underline{g}=\begin{cases}v^-\cos\bar\Delta & \bar\Delta\le\pi/2\\ v^+\cos\min(\bar\Delta,\pi) & \text{otherwise}\end{cases}\\
\bar d(\tau)&=\int_0^\tau v^+(s)\sin\min\big(\bar\Delta(s),\tfrac\pi2\big)\,ds, & \pi_k(\tau)&=2\ell_k\sin\big(\min(\Omega(\tau),\pi)/2\big).
\end{align*}
$\omega(v)$ is largest at $v^\star=\sqrt{a_\perp/\kappa_{\max}}$ (7.1\,m/s for $a_\perp=10$, $\kappa_{\max}=0.2$),
so $\bar\omega(s)=\kappa_{\max}v^\star$ if $v^\star\in[v^-,v^+]$, $\kappa_{\max}v^+$ if $v^+<v^\star$, and
$a_\perp/v^-$ if $v^->v^\star$.

\begin{proposition}[Physical occupancy of a vehicle]\label{prop:veh}
Under Assumption~\ref{ass:veh}, for every $\tau\in[0,T]$, $\Occ_k(\tau)\subseteq\Ohat_k^{\mathrm{phys}}(\tau)$ with
\[
L^-=S^--\pi_k-\ell'-\varepsilon_p,\qquad L^+=S^++\pi_k+\ell'+\varepsilon_p,\qquad
\Lambda=\bar d+\pi_k+w'+\varepsilon_p,
\]
where $\ell'=\ell'(\min(\bar\Delta,\pi/2))$ and $w'=w'(\min(\bar\Delta,\pi/2))$ from Lemma~\ref{lem:rot}.
\end{proposition}
\begin{proof}
By (V3), $v^-(s)\le v_k(s)\le v^+(s)$. By (V4), $|\psi_k(s)-\psi_k(0)|\le\Omega(s)$, hence by (V1)
$|\psi_k(s)-\hat\psi_k|\le\bar\Delta(s)$. By (V2) the rear axle moves by
$\int v_k(\cos,\sin)(\psi_k)\,ds$; its component along $u_k$ is $\int v_k\cos(\psi_k-\hat\psi_k)$, which lies
in $[S^-,S^+]$ (for the lower end use $v_k\cos\alpha\ge\underline g$, since $v_k\ge0$); its component along
$n_k$ is at most $\int v^+\sin\min(\bar\Delta,\pi/2)=\bar d$ in absolute value. By (V5) the centre moves
by the rear axle's displacement plus $l_k\big(h(\psi_k(\tau))-h(\psi_k(0))\big)$, whose norm is at most
$2\ell_k\sin(\Omega/2)=\pi_k$ (the chord identity again). The initial centre error adds $\varepsilon_p$.
Finally the box's heading differs from $\hat\psi_k$ by at most $\bar\Delta$, so by Lemma~\ref{lem:rot} it
is contained in the rectangle of half-dimensions $(\ell',w')$ about its centre.
\end{proof}

\paragraph{Reading the bound.} For small angles and constant speed $\hat v$: $S^\pm\approx(\hat v\pm\varepsilon_v)\tau\pm\tfrac12 a\tau^2$
(and $S^-$ stops growing once the object has stopped: it cannot reverse), and
$\bar d\approx\varepsilon_\psi\hat v\tau+\tfrac12 a_\perp\tau^2$ when $\bar\omega=a_\perp/\hat v$. A parked car
($\hat v=0$) has $S^-\approx0$: it cannot move backwards, and it moves sideways only by turning, which
it can only do while moving forwards ($\bar\omega=\kappa_{\max}v^+$).

\paragraph{What it costs.} The physical layer is what a car \emph{can} do, so next to moving traffic it
is very conservative: a car abreast in the next lane could, physically, turn into the ego within about
0.09\,s (scene E1). Physics alone cannot make $H$ useful in traffic. Its roles are to cap the calibrated
layer where data are thin (long $\tau$, rare situations), and to decide which objects are relevant at all
(Section~\ref{sec:guarantee}).

\subsection{Physical layer: pedestrians and other holonomic objects}
\begin{proposition}[Pedestrians]\label{prop:ped}
If $\|c_k(t)-\hat c_k\|\le\varepsilon_p$, $\|w_k(t)-\hat w_k\|\le\varepsilon_v$, $\|\dot w_k\|\le a_p$ and
$\|w_k\|\le v_{p,\max}$ over $[0,T]$, then
\[
\Occ_k(\tau)\subseteq\Big[\Disk\big(\hat c_k+\hat w_k\tau,\ \varepsilon_p+\varepsilon_v\tau+\tfrac12a_p\tau^2\big)\cap\Disk\big(\hat c_k,\ \varepsilon_p+v_{p,\max}\tau\big)\Big]\oplus\Disk_{r_k},
\]
with $r_k=\sqrt{\ell_k^2+w_k^2}$.
\end{proposition}
\begin{proof}
The first disk is the PDF's derivation (Taylor with integral remainder); the second follows from the
speed limit. The body lies within $r_k$ of its centre.
\end{proof}
Set-based pedestrian prediction \cite{koschi2018} uses $a_p=0.6$\,m/s$^2$ and $v_{p,\max}=2.0$\,m/s and
raises both when a pedestrian is observed to exceed them; in our data persons reach 2.62\,m/s and change
speed by up to 1.32\,m/s per second (p99.9), so we propose 3.0\,m/s and 1.5\,m/s$^2$.
The intersection of two disks is not a disk, but a sound test needs no more: the ego misses the
intersection whenever it misses either disk. So the margin against a pedestrian is the
\emph{largest} of the margins against the two disks.

\paragraph{Static objects} (class \texttt{static}) are their box enlarged by $\varepsilon_p$. Objects of an
unknown class are treated as pedestrians with vehicle-level speed and acceleration limits.

\subsection{Calibrated layer: what objects of this kind actually did}\label{sec:conformal}

\paragraph{Nominal prediction.} Any predictor can be used; the calibration below wraps whichever one is
chosen. We use constant velocity, projected on the heading for vehicles (they do not slide sideways):
the nominal box at $\tau$ is $\Rect(\hat c_k+\hat v_k\tau u_k,\ \hat\psi_k,\ \ell_k,\ w_k)$. For persons the
nominal centre is $\hat c_k+\hat w_k\tau$.

\paragraph{Scores.} For a calibration example $i$ (a known vehicle at some tick of a calibration clip),
let $q_{ij}(\tau)$, $j=1..4$, be the corners of its \emph{true} box at $t+\tau$, expressed in the frame
$(u,n)$ centred on the nominal centre. The \emph{overhangs} are how far the true box sticks out of the
nominal box:
\[
e^{f}_i(\tau)=\max_j q_{ij}\cdot u-\ell,\qquad e^{b}_i(\tau)=-\min_j q_{ij}\cdot u-\ell,\qquad e^{l}_i(\tau)=\max_j|q_{ij}\cdot n|-w,
\]
clipped at 0 (front, back, side). Rotation, turning and size changes are all included. With fixed
positive shape functions $g^a(\tau)$ (for instance the p90 curves from a separate fitting split), the
score of the example is its worst normalised overhang over the horizon:
\[
S_i=\max_{\tau\in\mathcal T_i}\ \max_{a\in\{f,b,l\}}\ e^a_i(\tau)/g^a(\tau),
\]
where $\mathcal T_i\subseteq\{0.1,0.2,\dots,6.4\}$ are the steps at which the object still exists.

\begin{theorem}[Split conformal prediction {\cite{vovk2005,angelopoulos2023}}]\label{thm:conformal}
Let the scores $S_1,\dots,S_n,S_{n+1}$ be exchangeable and let $\hat q$ be the
$\lceil(n+1)(1-\delta)\rceil$-th smallest of $S_1..S_n$ ($+\infty$ if that exceeds $n$). Then
$\Pr(S_{n+1}\le\hat q)\ge1-\delta$.
\end{theorem}
\begin{proof}[Proof sketch]
By exchangeability the rank of $S_{n+1}$ among the $n+1$ scores is uniform (ties only help), so it is
among the $\lceil(n+1)(1-\delta)\rceil$ smallest with probability at least $1-\delta$.
\cite{angelopoulos2023} state it for i.i.d.\ data and note that exchangeability suffices.
\end{proof}

\paragraph{The calibrated region.} With $\lambda^a(\tau)=\hat q\,g^a(\tau)$, the calibrated box of a new
object is the nominal box pushed out by $\lambda^f$ at the front, $\lambda^b$ at the back and $\lambda^l$ on
each side:
\[
L^-=\hat v_k\tau-\ell_k-\lambda^b(\tau),\qquad L^+=\hat v_k\tau+\ell_k+\lambda^f(\tau),\qquad \Lambda=w_k+\lambda^l(\tau).
\]
Because the score takes the maximum over $\tau$ and over the three sides, Theorem~\ref{thm:conformal} gives
\emph{one} probability for the whole horizon: with probability at least $1-\delta$ the true box stays in
the calibrated box at every step of $\mathcal T$. The alternative is to calibrate each step at level
$\delta/T$ and take a union bound: \cite{stankeviciute2021} do this for multi-step forecasts, and
\cite{lindemann2023} for planning against predicted agents (score $\|Y_\tau-\hat Y_\tau\|$, regions used as
constraints in a model-predictive controller). The maximum score avoids the $1/T$ loss when errors at
different steps are correlated, which they are here.

\paragraph{Groups.} Calibrate separately per group $G$ (class $\times$ stopped/moving, possibly speed
bins): the guarantee then holds within each group (Mondrian conformal prediction \cite{vovk2005}).
\cite{lindemann2023} make the same kind of assumption that is needed here: the ego's actions do not
change the distribution of the other agents' motion. The
data show why this matters (Table~\ref{tab:overhang}).

\begin{table}[h]
\centering
\begin{tabular}{rcccccc}
\toprule
& \multicolumn{3}{c}{stopped vehicles ($\|\hat w\|<0.5$\,m/s)} & \multicolumn{3}{c}{moving vehicles} \\
\cmidrule(lr){2-4}\cmidrule(lr){5-7}
$\tau$ [s] & front & back & side & front & back & side \\
\midrule
0.5 & 0.39 & 0.40 & 0.25 & 0.99 & 1.06 & 1.04 \\
1.0 & 0.75 & 0.76 & 0.43 & 2.21 & 2.24 & 2.02 \\
2.0 & 1.23 & 1.17 & 0.67 & 4.18 & 4.21 & 4.40 \\
4.0 & 2.16 & 1.95 & 0.86 & 8.76 & 9.45 & 10.72 \\
6.4 & 3.40 & 2.97 & 1.19 & 16.63 & 19.23 & 19.43 \\
\bottomrule
\end{tabular}
\caption{p99 of the box overhangs [m], per $\tau$, five clips, in-sample. The large side overhang of
moving vehicles comes mostly from cars turning at intersections: a car turning at 0.5\,rad/s swings its
front corner about 0.8\,m sideways in one second, on top of the curve itself.}
\label{tab:overhang}
\end{table}

\paragraph{How many examples.} $\hat q$ is finite only if $\lceil(n+1)(1-\delta)\rceil\le n$, i.e.
$n\ge(1-\delta)/\delta$: 99 examples per group for $\delta=1\%$, 999 for $\delta=0.1\%$. Examples from
the same track are strongly dependent, so each track should contribute one example (a random tick), and
calibration and evaluation clips must differ. Our five clips have 884 vehicle tracks and 192 VRU tracks
in total: enough for about $\delta=1\%$ per group today; $\delta=0.1\%$ needs roughly ten times more clips
(the catalog has them).

\paragraph{What the guarantee assumes.} Exchangeability: the new object is like the calibration
objects. A different city, weather or sensor stack breaks it; the safety check (T015, T016) on held-out
clips measures whether the coverage holds. The scores also inherit the labels' artefacts (Problem 5):
they are calibrated against the same labels that the safety check treats as truth, which is consistent.
Checking a miss rate needs many independent plans: with no failure in $n$ independent plans, the 95\%
upper bound on the rate is about $3/n$ \cite{hanley1983}.

\subsection{Combining the layers}\label{sec:combine}
If $\Occ_k\subseteq\Ohat_k^{\mathrm{phys}}$ always (Assumption~\ref{ass:veh} holds) and
$\Occ_k\subseteq\Ohat_k^{\mathrm{cal}}$ with probability $\ge1-\delta$, then
$\Occ_k\subseteq\Ohat_k=\Ohat_k^{\mathrm{phys}}\cap\Ohat_k^{\mathrm{cal}}$ with probability $\ge1-\delta$. Both are
rectangles on the axes $(u_k,n_k)$, so
\[
L_k^-=\max(L^{-,\mathrm{phys}},L^{-,\mathrm{cal}}),\qquad L_k^+=\min(L^{+,\mathrm{phys}},L^{+,\mathrm{cal}}),\qquad \Lambda_k=\min(\Lambda^{\mathrm{phys}},\Lambda^{\mathrm{cal}}).
\]
On our data the calibrated layer is the tighter one in all four example scenes; the physical layer matters
for long horizons, rare cases, and the relevance filter below.

%======================================================================
\section{Two horizons}\label{sec:horizons}

\subsection{Contact horizon}
\begin{definition}
With $\Ego(\tau)=\Rect\big(\text{footprint at }(\pref(\tau),\psiref(\tau))\big)$ enlarged by $\rho$,
\[
\Hc=\inf\Big\{\tau\in[0,\Tcap]:\ \exists k\in\mathcal K(t),\ \Ego(\tau)\cap\Ohat_k(\tau)\ne\emptyset\Big\},
\]
and $\Hc=\Tcap$ (capped) if there is no such $\tau$.
\end{definition}
This is the baseline's definition with the new sets. It answers ``how long can the plan be followed
without a possible contact''; it is the right quantity for question A of T023, but not a safe deadline
(Problem 4).

\subsection{Reaction horizon}
\begin{definition}[Fallback]
For a switch time $s\in[0,\Tcap]$, the fallback $F_s$ follows the plan until $s$ and then brakes along
the plan's path at $b_e$ until it stops. With $\sigma(\tau)$ the plan's arc length and
$v_s=\vref(s)$, its arc length is $\sigma(\tau)$ for $\tau\le s$ and
$\sigma(s)+v_s(\tau-s)-\tfrac12b_e(\tau-s)^2$ for $s<\tau\le\tau_{\mathrm{stop}}=s+v_s/b_e$; its pose is the
plan's pose at that arc length (extended straight beyond the plan's end). $\Ego_s(\tau)$ is its
footprint enlarged by $\rho$.
\end{definition}

\begin{definition}[Reaction horizon]
\[
\Hr=\sup\Big\{s\in[0,\Hc]:\ \forall s'\in[0,s],\ \forall\tau\in[0,\tau_{\mathrm{stop}}(s')],\ \forall k\in\mathcal K'(t):\ \Ego_{s'}(\tau)\cap\Ohat_k(\tau)=\emptyset\Big\},
\]
and $\Hr=0$ if even $F_0$ fails. $\mathcal K'(t)\subseteq\mathcal K(t)$ excludes the objects of
Assumption~\ref{ass:resp}.
\end{definition}

The car's deadline is then $t+\Hr-\Delta_{\mathrm{sw}}$, with $\Delta_{\mathrm{sw}}$ the time to take a
received plan into execution. Properties:
\begin{itemize}[leftmargin=*]
  \item $\Hr\le\Hc$, since $F_s$ coincides with the plan on $[0,s]$.
  \item $\Hr$ is a time-to-react of the plan: the longest time the plan can be continued before an
    evasive manoeuvre must, and still can, be executed \cite{hillenbrand2006,sontges2018}. Only braking
    along the path is tried, and the occupancies are over-approximations, so $\Hr$ is conservative: other
    fallbacks (swerving) can only make it larger, at more cost. Braking until rest as the safety
    certificate is the notion of passive safety of \cite{macek2009}.
  \item \textbf{Recursive safety.} If the next plan arrives at $t+s$ with $s\le\Hr$, the fallback $F_s$ of
    the \emph{current} plan is still clear. If the new plan's own $\Hr$ is 0, the car should not switch to
    it but execute $F_s$: it never has to accept an unverified plan. This is the scheme of online
    verification with fail-safe trajectories \cite{pek2020,pek2021}.
\end{itemize}

\begin{assumption}[Responsibility]\label{ass:resp}
\begin{enumerate}[label=(A\arabic*),leftmargin=*]
  \item \emph{Passive safety:} a contact that happens while the ego is at rest is not counted. Passive
    motion safety ``guarantees that, if a collision takes place, the robot will be at rest''
    \cite{bouraine2012}; \cite{macek2009} define it through a braking manoeuvre, as here. In $\Hr$, the
    check ends at $\tau_{\mathrm{stop}}$.
  \item \emph{Followers keep their distance:} a vehicle entirely behind the ego's rear bumper at $t$,
    heading within $45^\circ$ of the ego, is responsible for not running into it while the ego brakes at
    no more than $b_e$ \cite{rss2017}. Such vehicles are excluded from $\mathcal K'(t)$.
\end{enumerate}
\end{assumption}
Both are configuration switches. Without (A2), a fast car behind would make every braking fallback
unsafe and $\Hr=0$. The safety check (T015) must apply the same rules, and report rear-end contacts
separately: they are also the replay artefact T015 already anticipates (recorded objects do not react to
a plan that brakes differently from the real car).

\subsection{Car following reduces to RSS}
\begin{proposition}\label{prop:rss}
On a straight road, let the ego follow a plan at constant speed $v_e$ behind a vehicle with bumper gap
$g$ at $t$ and speed $\hat v_f$, $v_f^-=\hat v_f-\varepsilon_v\le v_e$. Use the physical layer without
heading change ($\Omega\equiv0$, $\varepsilon_\psi=0$) and $b_e\le b_{\max}$. Then, with
$m=\rho+\varepsilon_p$,
\[
\Hr=\frac{1}{v_e}\Big[g-m+\frac{(v_f^-)^2}{2b_{\max}}-\frac{v_e^2}{2b_e}\Big]_+ .
\]
\end{proposition}
\begin{proof}
The lead's rear is at least at $g+S^-(\tau)$, with $S^-$ the braking-to-stop distance from $v_f^-$ at
$b_{\max}$. The ego's front under $F_s$ is at $v_e\min(\tau,s)$ plus its braking distance after $s$. The
gap's derivative is the lead's speed minus the ego's; the lead starts no faster and brakes at least as
hard and earlier, so it is never faster and the gap shrinks until the ego stops. The worst instant is
therefore the ego's stop, where the condition reads $v_es+v_e^2/(2b_e)+m<g+(v_f^-)^2/(2b_{\max})$.
\end{proof}
This is the RSS minimum safe longitudinal distance \cite{rss2017} read the other way round: the response
time that the current gap still affords, with no acceleration during the response (the plan holds its
speed). For $g=4$\,m, $v_e=\hat v_f=4$\,m/s, $\varepsilon_v=1.5$, $\varepsilon_p=0.5$, $\rho=0.389$,
$b_{\max}=10$, $b_e=6$: $\Hr=(4-0.889+0.3125-1.333)/4=0.522$\,s; the numerical solver gives 0.52\,s on
its 0.01\,s grid, which checks the code. The same check for $\Hc$ (no fallback term): closed form
0.856\,s, numerical 0.86\,s.

\subsection{The guarantee}\label{sec:guarantee}
Call an object \emph{relevant} if its physical region can reach the region swept by the plan and its
fallbacks within the horizon; $\mathcal K_{\mathrm{rel}}\subseteq\mathcal K'(t)$. Objects that are not
relevant cannot cause a contact as long as the physical assumptions hold, whatever the statistics say.

\begin{theorem}\label{thm:guarantee}
Assume Lemma~\ref{lem:ego}'s tube, Assumptions~\ref{ass:veh}--\ref{ass:resp} (and the pedestrian
assumptions) and, for every $k\in\mathcal K_{\mathrm{rel}}$, calibrated regions with miss probability
$\delta_k$. Then the probability that the car, following the plan for any $s\le\Hr$ and then either a
new verified plan or $F_s$, touches a known object while moving is at most
$\sum_{k\in\mathcal K_{\mathrm{rel}}}\delta_k$.
\end{theorem}
\begin{proof}
On the event that every relevant object stays in its calibrated region, the definition of $\Hr$ makes
$F_s$ contact-free for every $s\le\Hr$, and a verified new plan is contact-free by its own $\Hr$. The
complement has probability at most $\sum_k\delta_k$ (union bound). Irrelevant objects cannot reach the
ego by the physical layer.
\end{proof}
To promise a per-plan probability $\delta_{\mathrm{plan}}$, split it: $\delta_k=\delta_{\mathrm{plan}}/|\mathcal K_{\mathrm{rel}}|$.
Objects \emph{unknown} at $t$ (not yet seen, or not observed) are outside the guarantee; T015 counts their
contacts separately, as decided in T023.

%======================================================================
\section{Computation}\label{sec:computation}

\paragraph{Time grid.} Both horizons are found on a grid of step $\Delta$ ($0.01$\,s), then refined by
bisection. Between grid points the true bodies move by at most half a step at their maximum point speed,
so a continuous-time guarantee follows by enlarging each set by $\tfrac12\Delta$ times that speed
(ego: $\vref+\dot\psi_{\mathrm{ref}}R_F$; object: $v^++\bar\omega\sqrt{\ell^2+w^2}$): 7.5\,cm at 15\,m/s, or
choose a finer grid.

\paragraph{Algorithm, per plan.}
\begin{enumerate}[leftmargin=*]
  \item Known objects $\mathcal K(t)$; drop those excluded by (A2); for each, precompute the physical
    profile ($v^\pm,S^\pm,\bar d,\Omega$: cumulative sums on the grid) and look up its group's calibrated
    curves $\lambda^a(\tau)$.
  \item Relevance filter with the physical layer (cheap: one bounding test per object).
  \item $\Hc$: first grid $\tau$ at which $\Ego(\tau)$ meets some $\Ohat_k(\tau)$ (rectangle--rectangle or
    rectangle--disk), then bisection.
  \item $\Hr$: for $s$ on a coarser grid (e.g.\ 0.05\,s) up to $\Hc$, check $F_s$ over
    $[s,\tau_{\mathrm{stop}}(s)]$ (the part before $s$ is the plan, already clear); $\Hr$ is the last $s$
    before the first failure, refined by bisection (binary search is also how \cite{tamke2011} compute
    such time measures quickly with many objects).
\end{enumerate}
Cost: $O(K\cdot N_\tau)$ constant-time tests for $\Hc$ and $O(K\cdot N_s\cdot N_{\mathrm{stop}})$ for $\Hr$;
with $K\approx30$, $N_s=128$ and $N_{\mathrm{stop}}\le300$ this is about $10^6$ tests per plan, which
vectorises.

%======================================================================
\section{Worked examples}\label{sec:examples}
All numbers come from \texttt{scripts/t024\_worked\_examples.py}. Common settings: plan straight along
$+x$; ego footprint as above, $\rho=0.389$\,m; physical layer with the values of Table~\ref{tab:params};
calibrated layer with the provisional p99 overhang curves (Table~\ref{tab:params}); baseline with the
repository's configuration and $r_{\mathrm{veh}}=3.79$\,m; median car $3.89\times1.77$\,m.

\paragraph{E1: car in the next lane, abreast, both at 10\,m/s.} Centres 3.5\,m apart sideways; the real
gap between the boxes is $3.5-1.0-0.885=1.615$\,m. Baseline: the disks need 6.43\,m at $\tau=0$, so
$H=0$. Revised: the side room is $1.615-0.389=1.226$\,m. Physically a car at about 9\,m/s may turn at
$a_\perp/v\approx1.1$\,rad/s, and its box swings out by $\pi_k+w'-w$; this eats the room in 0.09\,s. The
calibrated side overhang of moving cars, $1.94\tau+0.175\tau^2$, reaches 1.226\,m at 0.60\,s. So
$\Hc=0.60$\,s: short, because moving cars in our data often turn; a lane map, or conditioning on the
object's own yaw rate, is what would lengthen it.

\paragraph{E2: following at 4\,m/s, 4\,m gap.} Baseline: centre distance $3.656+4+1.945=9.60$\,m against
$6.43+1.0\tau+\tau^2$, so $H=1.36$\,s. Revised, calibrated: the lead's back overhang
$0.74+1.02\tau+0.29\tau^2$ reaches the room $4-0.389$ at $\Hc=1.85$\,s; braking at 6\,m/s$^2$ after
$s$ takes $0.67$\,s and 1.33\,m, which leaves $\Hr=1.66$\,s. Physical: the lead may brake at
10\,m/s$^2$ and turn, $\Hc=0.50$\,s, $\Hr=0.03$\,s; without turning, $\Hc=0.86$\,s and $\Hr=0.52$\,s, as
Proposition~\ref{prop:rss} predicts.

\paragraph{E3: the pedestrian of the hand example.} Pedestrian at $(30,-8)$\,m walking at 1.5\,m/s
towards the path; plan at 10\,m/s. Baseline 2.00\,s (2.04\,s in \texttt{tests/hand\_example.md}, which
uses $r_{\mathrm{veh}}=2$ and $\bar e_\theta=0.1$). Revised: physical 2.23\,s; calibrated (person p99,
$0.85\tau+0.025\tau^2$ plus the body radius 0.42\,m) 2.57\,s. Here the change is small, because the
pedestrian really is in the way: boxes help most beside the path, not on it.

\paragraph{E4: passing a parked car.} Parked car 15\,m ahead, centre 3.0\,m to the side (gap
$3.0-1.0-0.885=1.115$\,m), ego at 10\,m/s. Baseline 0.76\,s, physical 0.75\,s (physically a parked car can
pull out). Calibrated: stopped cars' side overhang is $0.28+0.15\tau$, which never reaches the room
$1.115-0.389=0.726$\,m before the ego has passed: capped at 6.4\,s.

%======================================================================
\section{Parameters}\label{sec:params}

\begin{table}[h]
\centering
\small
\begin{tabular}{p{2.3cm}p{3.3cm}p{4.4cm}p{5.4cm}}
\toprule
Symbol & Meaning & Proposed & Source \\
\midrule
$\rho_p$ & ego tracking tube & 0.2\,m & assumption (was $\rho_0$); controller spec, T018 \\
$\bar e_\theta$ & ego heading error & 0.05\,rad & assumption (current config) \\
$R_F$ & farthest footprint point & 3.790\,m & data (footprint) \\
$b_e$ & ego fallback braking & 6\,m/s$^2$ & assumption: above 3.4 (design braking \cite{aashto2001}), below ABS max \\
$\Delta_{\mathrm{sw}}$ & plan switch latency & 0.1\,s & assumption \\
\midrule
$\varepsilon_p$ & position error (phys.) & 0.5\,m & data: parked RMS p50 0.16, p95 0.35\,m \\
$\varepsilon_\psi$ & heading error (phys.) & $6^\circ$ & data: parked yaw RMS p50 $1.0^\circ$, p95 $2.8^\circ$ \\
$\varepsilon_v$ & speed error (phys.) & 1.5\,m/s & data: T009 p95 0.47--0.72\,m/s; tails larger \\
$a_{\max}$ & max acceleration & 4\,m/s$^2$ & data: 1\,s windows max 4.6\,m/s$^2$ \\
$b_{\max}$ & max braking & 10\,m/s$^2$ & ABS, dry road: 8.5--10.6 \cite{miller2023}; data min $-7.2$ \\
$v_{\max}$ & max speed & 50\,m/s & data: max 45.9\,m/s \\
$\kappa_{\max}$ & max curvature & 0.2\,m$^{-1}$ & assumption (5\,m radius); data p99 0.11 \\
$a_\perp$ & max lateral acceleration & 10\,m/s$^2$ & friction as for $b_{\max}$; normal driving $\le4$ \cite{bosetti2014} \\
$a_p,\ v_{p,\max}$ & pedestrian limits & 1.5\,m/s$^2$, 3\,m/s & data: p99.9 1.3, max 2.6; \cite{koschi2018,bohannon1997} \\
\midrule
$\lambda^{f,b,l}$ stopped & calibrated, p99 & $0.21{+}0.49\tau$, $0.28{+}0.42\tau$, $0.28{+}0.15\tau$ & data, in-sample, provisional \\
$\lambda^{f,b,l}$ moving & calibrated, p99 & $0.48{+}1.39\tau{+}0.18\tau^2$, $0.74{+}1.02\tau{+}0.29\tau^2$, $1.94\tau{+}0.18\tau^2$ & data, in-sample, provisional \\
$\lambda$ person & calibrated, p99 & $0.85\tau+0.025\tau^2$ (centre) & data, in-sample, provisional \\
$\delta$ & miss probability per object & 1\% & limited by sample size (Section~\ref{sec:conformal}) \\
\bottomrule
\end{tabular}
\caption{Proposed parameters. The physical-layer error bounds are set well above the measured p95,
because a physical bound must hold almost always; the calibrated layer carries the realistic tails.
UN Regulation 157 likewise requires an automated lane-keeping system to avoid a lead vehicle braking
``up to its full braking performance'' \cite{unr157}.}
\label{tab:params}
\end{table}

%======================================================================
\section{What this changes in the pipeline (proposal)}\label{sec:pipeline}
Nothing below is implemented yet; each item would be its own ticket.
\begin{enumerate}[leftmargin=*]
  \item \textbf{New solver \texttt{method: boxes}} in the runner's registry: Lemmas
    \ref{lem:ego}--\ref{lem:rot}, the calibrated layer and $\Hc$. This alone removes most $H=0$
    plans (Problem 1).
  \item \textbf{Calibration experiment} (T019 rewritten): the conformal scores of
    Section~\ref{sec:conformal} per group, on calibration clips, evaluated on held-out clips; more clips
    from the catalog for $\delta\le0.1\%$.
  \item \textbf{Reaction horizon} $\Hr$ in the solver, and an amendment to T023: the deadline uses $\Hr$.
  \item \textbf{Safety check} (T015): the same boxes; besides ``contact before $H$'', check the fallback
    $F_H$ against the true future; label rear-end contacts and contacts at rest (Assumption~\ref{ass:resp}).
  \item \textbf{Physical layer} (Propositions~\ref{prop:veh}--\ref{prop:ped}) as relevance filter and cap.
\end{enumerate}

\paragraph{Decisions for Ralph.}
\begin{enumerate}[label=D\arabic*.,leftmargin=*]
  \item The miss probability: $\delta$ per object, or $\delta_{\mathrm{plan}}$ split over relevant objects.
  \item The deadline: $\Hc$ or $\Hr$ (recommended: $\Hr$, Problem 4).
  \item The fallback braking $b_e$ and the switch latency $\Delta_{\mathrm{sw}}$.
  \item The responsibility rules (A1), (A2): on or off.
  \item The ego tube $\rho_p$, $\bar e_\theta$ (no data until a controller exists).
  \item The nominal predictor for vehicles: heading-projected constant velocity, or constant turn rate
    (would shrink the side overhang of turning cars, if the yaw rate can be estimated causally; T009 found
    single-step object yaw rates noisy).
\end{enumerate}

%======================================================================
\section{Limitations}\label{sec:limits}
\begin{itemize}[leftmargin=*]
  \item \textbf{Labels as truth.} Calibration and the safety check both use the auto-labels; their
    artefacts (jumps of several metres) inflate the tails. A real perception stack has different errors.
  \item \textbf{Objects do not react} in the replay; a plan that brakes harder than the real car can be
    ``rear-ended'' (handled by (A2) and reported separately).
  \item \textbf{No map.} Lane geometry is the strongest constraint on other vehicles: set-based
    prediction keeps them on the road and in lanes they may legally use (constraint C5 of
    \cite{althoff2016}, implemented in SPOT \cite{koschi2017}). Without it, the lateral bounds of moving
    cars are wide (E1).
  \item \textbf{Unknown objects} (not yet seen, occluded) are outside the guarantee by construction.
  \item \textbf{Small data.} Five clips, about 900 vehicle tracks: $\delta\approx1\%$ per group at best;
    the calibrated numbers in this draft are in-sample and per-$\tau$.
  \item \textbf{Censoring.} Tracks that leave the scene early contribute only short horizons; if fast or
    erratic objects leave sooner, long-$\tau$ quantiles are biased low.
\end{itemize}

%======================================================================
\appendix
\section{How the data numbers were computed}\label{app:data}
\texttt{scripts/spike\_t024\_object\_motion.py}, offline scenes of the five clips, objects with
\texttt{in\_scene} at $t$ and a known velocity; futures from the same track at tick $t+10\tau$.
\begin{itemize}[leftmargin=*]
  \item \textbf{Residuals:} $r=c(t+\tau)-\hat c-\hat w\tau$ (full velocity) or $-\hat v u\tau$
    (heading-projected), split along $u$ and $n$ at $t$.
  \item \textbf{Overhangs:} the true box's corners at $t+\tau$ against the heading-projected nominal box,
    as defined in Section~\ref{sec:conformal}.
  \item \textbf{1\,s chords} (motion limits, robust to label noise): $c_1=c(t+1)-c(t)$,
    $c_2=c(t+2)-c(t+1)$; longitudinal acceleration $\|c_2\|-\|c_1\|$, heading change $\angle(c_1,c_2)$,
    lateral acceleration $\bar v\angle(c_1,c_2)$, curvature $\angle/\bar v$. One-second averaging lowers
    peaks, so these are lower bounds on the true extremes.
  \item \textbf{Label noise:} tracks that look parked (T009's rule): RMS distance from their mean position
    and RMS yaw about their mean.
  \item \textbf{$H=0$ plans:} from \texttt{results/t014-baseline} (T014), the ego footprint and the
    limiting object's box at the plan's tick; exact rectangle distance.
\end{itemize}

\begin{thebibliography}{99}\small
\bibitem{rss2017} S.~Shalev-Shwartz, S.~Shammah, A.~Shashua. On a formal model of safe and scalable self-driving cars. arXiv:1708.06374v6, 2017. (Minimal safe longitudinal distance: Lemma~2.)
\bibitem{althoff2016} M.~Althoff, S.~Magdici. Set-based prediction of traffic participants on arbitrary road networks. \emph{IEEE Transactions on Intelligent Vehicles}, 1(2):187--202, 2016. doi:10.1109/TIV.2016.2622920.
\bibitem{koschi2017} M.~Koschi, M.~Althoff. SPOT: a tool for set-based prediction of traffic participants. \emph{IEEE Intelligent Vehicles Symposium}, 1686--1693, 2017. doi:10.1109/IVS.2017.7995951.
\bibitem{koschi2018} M.~Koschi, C.~Pek, M.~Beikirch, M.~Althoff. Set-based prediction of pedestrians in urban environments considering formalized traffic rules. \emph{IEEE ITSC}, 2704--2711, 2018. doi:10.1109/ITSC.2018.8569434.
\bibitem{pek2020} C.~Pek, S.~Manzinger, M.~Koschi, M.~Althoff. Using online verification to prevent autonomous vehicles from causing accidents. \emph{Nature Machine Intelligence}, 2(9):518--528, 2020. doi:10.1038/s42256-020-0225-y.
\bibitem{pek2021} C.~Pek, M.~Althoff. Fail-safe motion planning for online verification of autonomous vehicles using convex optimization. \emph{IEEE Transactions on Robotics}, 37(3):798--814, 2021. doi:10.1109/TRO.2020.3036624.
\bibitem{sontges2018} S.~S\"ontges, M.~Koschi, M.~Althoff. Worst-case analysis of the time-to-react using reachable sets. \emph{IEEE Intelligent Vehicles Symposium}, 1891--1897, 2018. doi:10.1109/IVS.2018.8500709.
\bibitem{hillenbrand2006} J.~Hillenbrand, A.~M.~Spieker, K.~Kroschel. A multilevel collision mitigation approach---its situation assessment, decision making, and performance tradeoffs. \emph{IEEE Transactions on Intelligent Transportation Systems}, 7(4):528--540, 2006. doi:10.1109/TITS.2006.883115. (Time-to-react, as credited by \cite{sontges2018}.)
\bibitem{tamke2011} A.~Tamke, T.~Dang, G.~Breuel. A flexible method for criticality assessment in driver assistance systems. \emph{IEEE Intelligent Vehicles Symposium}, 697--702, 2011. doi:10.1109/IVS.2011.5940482.
\bibitem{macek2009} K.~Ma\v{c}ek, D.~Vasquez, T.~Fraichard, R.~Siegwart. Towards safe vehicle navigation in dynamic urban scenarios. \emph{Automatika}, 50(3--4):184--194, 2009.
\bibitem{bouraine2012} S.~Bouraine, T.~Fraichard, H.~Salhi. Provably safe navigation for mobile robots with limited field-of-views in dynamic environments. \emph{Autonomous Robots}, 32(3):267--283, 2012. doi:10.1007/s10514-011-9258-8.
\bibitem{ziegler2010} J.~Ziegler, C.~Stiller. Fast collision checking for intelligent vehicle motion planning. \emph{IEEE Intelligent Vehicles Symposium}, 518--522, 2010. doi:10.1109/IVS.2010.5547976.
\bibitem{tolksdorf2024} L.~Tolksdorf, C.~Birkner, A.~Tejada, N.~van~de~Wouw. Fast collision probability estimation for automated driving using multi-circular shape approximations. \emph{IEEE Intelligent Vehicles Symposium}, 2024. arXiv:2405.10765.
\bibitem{gilbert1988} E.~G.~Gilbert, D.~W.~Johnson, S.~S.~Keerthi. A fast procedure for computing the distance between complex objects in three-dimensional space. \emph{IEEE Journal on Robotics and Automation}, 4(2):193--203, 1988. doi:10.1109/56.2083.
\bibitem{gottschalk1996} S.~Gottschalk, M.~C.~Lin, D.~Manocha. OBBTree: a hierarchical structure for rapid interference detection. \emph{SIGGRAPH '96}, 171--180, 1996. doi:10.1145/237170.237244.
\bibitem{mayne2005} D.~Q.~Mayne, M.~M.~Seron, S.~V.~Rakovi\'c. Robust model predictive control of constrained linear systems with bounded disturbances. \emph{Automatica}, 41(2):219--224, 2005. doi:10.1016/j.automatica.2004.08.019.
\bibitem{vovk2005} V.~Vovk, A.~Gammerman, G.~Shafer. \emph{Algorithmic Learning in a Random World}. Springer, 2005. doi:10.1007/b106715.
\bibitem{angelopoulos2023} A.~N.~Angelopoulos, S.~Bates. Conformal prediction: a gentle introduction. \emph{Foundations and Trends in Machine Learning}, 16(4):494--591, 2023. doi:10.1561/2200000101. Preprint arXiv:2107.07511.
\bibitem{stankeviciute2021} K.~Stankevi\v{c}i\={u}t\.{e}, A.~M.~Alaa, M.~van~der~Schaar. Conformal time-series forecasting. \emph{Advances in Neural Information Processing Systems}, 34:6216--6228, 2021.
\bibitem{lindemann2023} L.~Lindemann, M.~Cleaveland, G.~Shim, G.~J.~Pappas. Safe planning in dynamic environments using conformal prediction. \emph{IEEE Robotics and Automation Letters}, 8(8):5116--5123, 2023. doi:10.1109/LRA.2023.3292071.
\bibitem{hanley1983} J.~A.~Hanley, A.~Lippman-Hand. If nothing goes wrong, is everything all right? Interpreting zero numerators. \emph{JAMA}, 249(13):1743--1745, 1983. doi:10.1001/jama.1983.03330370053031.
\bibitem{miller2023} I.~Miller, D.~King, C.~Wilkinson, G.~P.~Siegmund. Decelerations for vehicles with anti-lock brake systems (ABS) on dry asphalt and concrete road surfaces. SAE Technical Paper 2023-01-0616, 2023. doi:10.4271/2023-01-0616.
\bibitem{aashto2001} AASHTO. \emph{A Policy on Geometric Design of Highways and Streets}. Washington, DC, 2001.
\bibitem{bosetti2014} P.~Bosetti, M.~Da~Lio, A.~Saroldi. On the human control of vehicles: an experimental study of acceleration. \emph{European Transport Research Review}, 6(2):157--170, 2014. doi:10.1007/s12544-013-0120-2.
\bibitem{bohannon1997} R.~W.~Bohannon. Comfortable and maximum walking speed of adults aged 20--79 years: reference values and determinants. \emph{Age and Ageing}, 26(1):15--19, 1997. doi:10.1093/ageing/26.1.15.
\bibitem{unr157} UNECE. UN Regulation No.~157 --- Automated Lane Keeping Systems. E/ECE/TRANS/505/Rev.3/Add.156/Amend.1, 2022.
\end{thebibliography}

\end{document}
